\section{Speedup over Exhaustive Search}\label{sec:advantage}

This section states the speedup that the interface of Section~\ref{sec:interface} provides and delimits its scope. The comparison is made in the \emph{query model}, in which a classical algorithm accesses the objective function only by evaluating $F(x)$ at plans of its choice and a quantum algorithm accesses it only through the oracles $\Otau$.

\begin{proposition}\label{prop:advantage}
Let $G$ be an instance of the cardinality variant of the MFNIP with $n$ nodes, $m$ arcs, and budget $k$, and let $\Fk=\{x\in\{0,1\}^m:|x|\le k\}$ and $N=|\Fk|$.
\begin{enumerate}
\item[(a)] Algorithm~\ref{alg:dh}, with the preparation $\Ak$ and the oracles $\Otau$ each implemented with $\poly(n,m)$ gates, returns an optimal interdiction plan with probability at least $1/2$ when run with query budget $Q=22.5\sqrt{N}+1.4\log_2^2 N$. It therefore uses $\bigO(\sqrt{N})$ applications of $\Ak$ and $\Otau$ and $\bigO(\sqrt{N}\cdot\poly(n,m))$ gates.
\item[(b)] Every classical algorithm that accesses $F$ only through evaluations and returns an optimal plan with probability at least $1/2$ uses at least $N/2-1$ evaluations on some instance of the black-box problem, and exhaustive enumeration of $\Fk$ uses $N$.
\item[(c)] Every quantum algorithm that accesses $F$ only through these oracles needs $\Omega(\sqrt{N})$ queries, so in this model the bound in (a) is optimal up to a constant factor.
\end{enumerate}
\end{proposition}

\begin{proof}
\emph{Efficient preparation.} The operator $\Ak$ consists of one rotation, $k-1$ controlled rotations, and the Dicke unitary $U_{m,k}$, which has $\bigO(km)$ gates~\cite{bartschi2019}. Since $k\le m$, the total is $\bigO(m^2)$ gates. Because $U_{m,k}$ maps $\ket{0^{m-\ell}1^\ell}$ to $\ket{D^m_\ell}$ for every $\ell\le k$, the state $\Ak\ket{0^m}$ is the uniform superposition over $\Fk$, as shown in Section~\ref{sec:interface-space}.

\emph{Efficient oracle.} A maximum $s$--$t$ flow can be computed classically in time polynomial in $n$ and $m$, for example with $\bigO(nm^2)$ arithmetic operations on numbers of $\bigO(\log(mU))$ bits, where $U=\max_a u_a$~\cite{ahuja1993}. By Bennett's construction~\cite{bennett1973}, a classical computation of time $T$ can be simulated reversibly in time $\bigO(T)$ and space $\bigO(T)$ by computing while keeping a history, copying the output, and uncomputing, which leaves only the input and the output. Hence the map $\ket{x}\ket{0}\mapsto\ket{x}\ket{F(x)}$ has a circuit with $\poly(n,m)$ gates. A reversible comparison with $\tau$ recorded in a flag qubit, a $Z$ gate on that qubit, and the inverse of the circuit form $\Otau$ with $\poly(n,m)$ gates. Removing the interdicted arcs only multiplies each capacity $u_a$ by $1-x_a$, which is an operation controlled by the input.

\emph{Amplification from $\Ak$.} Let $M_\tau=|\{x\in\Fk : F(x)<\tau\}|$ and $\sin^2\theta=M_\tau/N$. With $G=\Ak S_0\Ak^{\dagger}\Otau$, the state after $j$ iterations stays in the plane spanned by the normalized uniform superpositions of the marked and unmarked plans of $\Fk$, and a marked plan is measured with probability $\sin^2((2j+1)\theta)$~\cite{bhmt2002}. Conditioned on being marked, the measured plan is uniformly distributed over the marked set. This is exactly the behavior of Grover's algorithm on an unsorted table of $N$ items, so the analyses of exponential search~\cite{bbht1998} and of minimum finding~\cite{durr1996}, which rely only on this behavior, hold with the table replaced by $\Fk$. Dürr and Høyer charge $\log_2 N$ steps for preparing the initial state and one step per Grover iteration, and they show that the expected time until the threshold equals the minimum is at most $Q/2$~\cite{durr1996}. By Markov's inequality, a run interrupted at time $Q$ succeeds with probability at least $1/2$. Every measurement, and hence every classical evaluation, follows a preparation, so for $N\ge 2$ our cost measure (queries plus evaluations, Section~\ref{sec:interface-cost}) never exceeds their time count, and the same budget $Q$ bounds it. Each Grover iteration applies $\Otau$ once and $\Ak$ and $\Ak^{\dagger}$ once each, and each round applies $\Ak$ once more, so $\Ak$ and $\Otau$ are applied at most $3Q=\bigO(\sqrt{N})$ times. Multiplying this number by the gate cost of these operators gives the gate bound in (a).

\emph{Classical lower bound.} Consider black-box objectives $f:\Fk\to\{0,1\}$ with a single zero at an unknown position $x^\star$ and value one on every other plan. With $x^\star$ uniformly distributed, a randomized algorithm that makes $q$ evaluations and then returns a plan succeeds with probability at most $(q+1)/N$, since it either evaluated $x^\star$ or must guess it among the plans it did not evaluate. Success probability $1/2$ therefore requires $q\ge N/2-1$, and by Yao's principle the same bound holds for the worst case of any randomized algorithm. Exhaustive enumeration evaluates all $N$ plans and certifies the minimum.

\emph{Quantum lower bound.} The same family of objectives reduces minimum finding to unstructured search over $N$ items, which requires $\Omega(\sqrt{N})$ oracle queries for constant success probability~\cite{bbbv1997,zalka1999}.
\end{proof}

\begin{remark}[Scope of the proposition]
The speedup in Proposition~\ref{prop:advantage} is quadratic and asymptotic, it holds in the query model, and it is relative to exhaustive search over $\Fk$. It says nothing about classical algorithms that exploit the structure of the network. The proposition matters only when $N$ is large, that is, when the budget grows with the network. For $k\le m/2$ we have $N\ge\binom{m}{k}\ge(m/k)^k$, which is superpolynomial in $m$ whenever $k$ grows with $m$ (for example, $k=\Theta(\sqrt{m})$ or $k=\Theta(m)$). For constant $k$, $N=\bigO(m^k)$ and both searches are polynomial.
\end{remark}
